\section{Discussion}\label{sec:discussion}

\subsection{Why the particle filter wins, and where the EKF breaks}
The decisive quantity is the ratio between the position uncertainty and the distance over which
the terrain slope stays roughly constant. The EKF replaces the terrain by its tangent plane at the
current estimate; that is adequate only when the error is small compared with the terrain's
correlation length. The smooth synthetic hills used during development have correlation lengths of
hundreds of metres to kilometres, and there the EKF held up to $\sigma_0 = 1$\,km. The real 30\,m
DEM of the hill country changes slope within a few cells, so the EKF already fails in
\EOneHillSOneZeroZeroEkfFl\,\% of the runs at $\sigma_0 = 100$\,m (E1). E2 shows the same effect
from the other side: the rougher the terrain under a flight, the \emph{larger} the EKF's error
($\rho = +\ETwoRhoEkf$) and the smaller the particle filter's ($\rho = \ETwoRhoPfFiveZeroZeroZero$).
Terrain texture is information for an estimator that evaluates the likelihood globally and noise
for one that linearises it locally.

Two consequences are worth stressing. First, a failed EKF is not merely inaccurate: its error
grows faster than that of the uncorrected INS, because corrections computed from the wrong slope
also corrupt its bias estimate (\cref{fig:e1time}). A linearised TAN filter whose linearisation is
invalid is worse than no terrain aiding. Second, its failures are confident: the EKF's dominant
failure mode is false lock (\cref{fig:e1rates}), so its own covariance gives no warning.

The particle filter pays for its robustness in computation, but on this problem the cost is
irrelevant: \RtPfFiveZeroZeroZero\,ms per step for 5\,000 particles on one laptop core, three
orders of magnitude inside the 1\,Hz measurement interval (E3).

\subsection{Terrain decides}
The three regions differ by two orders of magnitude in the information one altimeter sample
carries (\cref{tab:terrain}), and that ordering carries through every experiment. Over the hills
the particle filter tracks at 6--10\,m; over the plains at about 20--25\,m, close to the posterior
bound there (E5). The plains also show the most desirable property of a Bayesian filter: when the
terrain says little, the particle filter's covariance stays wide and it almost never locks
falsely (E1), so a navigation system could trust its own uncertainty and fall back on other
sources. The EKF has no such honesty --- it false-locks most often precisely where the terrain is
least informative relative to its linearisation errors.

The along-/cross-track decomposition (E2) did not show the elongated error along ridges and valleys
one might expect: averaged over flights with a 25$^\circ$ heading weave over 40\,km of varied
terrain, both directions are about equally constrained. The effect may exist on individual straight
flights along a valley, which this design does not isolate.

\subsection{Failure modes and their remedies}
\paragraph{False lock.} With a wide prior the particle filter's dominant failure is to collapse onto
the wrong patch of terrain (E1, E3). Three remedies were tested on identical flights: more particles
(E3: \EThreeNFiveKConv\,\% $\to$ \EThreeNTwentyKConv\,\% convergence from 5\,000 to 20\,000 at
$\sigma_0 = 1$\,km), roughening (E3: \EThreeRoughZeroPointZeroFiveConv\,\% with $K = 0.05$ at the
cost of 5\,000 particles) and the divergence monitor (E6). Roughening acts early, by keeping
hypotheses near the truth alive while the cloud is wide; the monitor acts late, by detecting a lock
that has already happened and re-opening the search.
On the main seeds (E6) the monitor alone lifts convergence at $\sigma_0 = 2$\,km on the hill
country from \ESixHillSTwoZeroZeroZeroPfFiveZeroZeroZeroConv\,\% to
\ESixHillSTwoZeroZeroZeroPfMonitorConv\,\%, roughening alone to
\ESixHillSTwoZeroZeroZeroPfRoughConv\,\%, and the two together to
\ESixHillSTwoZeroZeroZeroPfMonitorRoughConv\,\% with almost no false locks --- more than four
times as many particles achieve without them. The monitor's value lies in a property no amount of
particles provides: it turns a silent failure into a detected one, from information available on
board (the predicted innovation), without ground truth. Its false alarms are cheap because the
re-spread keeps half of the current cloud. Where the filter fails by imprecision rather than by
locking (the plains), there is nothing for the monitor to detect, and only roughening helps.

\paragraph{Altimeter outliers.} A Gaussian likelihood lets a single gross radar error destroy the
particles at the truth. The Student-$t$ likelihood keeps convergence at
\EFourOutliersHillOneZeroPfStudentTConv\,\% with 10\,\% outliers (Gaussian:
\EFourOutliersHillOneZeroPfFiveZeroZeroZeroConv\,\%) and costs nothing without them (E4). There is
no reason not to use it.

\paragraph{Map and altimeter offsets.} A random map error of the size the Copernicus handbook
implies is harmless, but a constant offset between measured and mapped height --- a barometric
bias, or a radar seeing through a canopy the surface model includes --- defeats the filter beyond
about 5\,m (E4). Inflating $\sigma_z$ buys tolerance up to about 10\,m; estimating the offset as a
state is the only remedy at 20\,m but costs acquisition when there is no offset. The obvious next
step is to treat the offset like the bias in the RBPF: it enters the measurement linearly, so it
could be marginalised analytically instead of sampled.

\paragraph{Static-parameter impoverishment.} The plain particle filter cannot estimate the INS bias
(E7): with a nearly static bias, resampling leaves bias values from a few ancestors, and the
estimate is no better than the prior. Rao-Blackwellisation removes the problem at no extra cost and
improves tracking from 2\,000 particles. It is the better filter for this model.

\subsection{How close to optimal?}
The posterior Cramér--Rao bound is never violated, which checks the whole chain of simulation,
filters and bound. In the tracking phase the filters are about twice the per-flight bound on the
hills, and the mismatched-gain analysis shows why: the deliberately inflated process noise
($\sigma_p = 2$\,m per step vs.\ a true $\approx0.05$\,m) alone allows only
\EFivePlanHillPfTwoZeroZeroZeroZeroTuned\,m, and the 20\,000-particle PF reaches
\EFivePlanHillPfTwoZeroZeroZeroZeroRmse\,m (E5). The inflation is a deliberate trade:
smaller process noise would track closer to the bound but would make the cloud less diverse during
acquisition, where E3 shows diversity (roughening) to be what matters. During acquisition the
bound is not a useful benchmark, because it ignores ambiguity --- which is precisely what makes TAN
hard; a bound that accounts for ambiguity would be needed there. The NEES behaves accordingly: consistent to conservative in the tracking phase,
uninformative while the posterior is multimodal.

\subsection{Limitations}\label{sec:limitations}
\begin{itemize}
  \item \textbf{Simulated sensors and flights.} No real flight data were used. The altimeter noise
        (2\,m barometric, 3\,m radar, white and range-independent) is an assumption, and an
        optimistic one at the 1--1.7\,km radar ranges the corridor altitude produces over the
        hills (E1).
  \item \textbf{Simplified flight.} Constant altitude and a horizontal 2-D state; attitude, height
        and their errors are not estimated, and the INS model is a velocity bias plus white noise.
  \item \textbf{The world uses the same map.} The simulated radar sees the DEM itself unless E4
        adds a map error; real canopy and DSM effects are spatially structured and only
        approximated by a smooth random field and a constant offset.
  \item \textbf{One tuning.} Process and measurement noise were fixed, not optimised per filter;
        the EKF in particular might do somewhat better with different tuning, although its failures
        are structural (linearisation), not a matter of noise levels.
  \item \textbf{Three regions of one country.} The conclusions about terrain dependence rest on
        three 40\,km boxes; they are consistent with the per-flight analysis of E2, but other
        landscapes (urban, forest, desert dunes) were not tested.
\end{itemize}

\subsection{Future work}
Natural extensions are a Rao-Blackwellised offset state, an adaptive proposal that uses the
current measurement (auxiliary or likelihood-based proposals) to fight false lock at its source,
a point-mass filter as a deterministic reference, real flight data with a recorded radar
altimeter, and an ambiguity-aware performance bound for the acquisition phase.
